\section{Clearance coarea and its rank-deficient positive remainder}
\label{sec:v61-clearance-coarea}
A small competing-hit clearance does not supply the incidence factor
in the preceding determinant. Its density instead admits a direct
estimate where the total roof and the image-side clearance coordinate
are transverse. We prove that estimate on every original physical word
and retain the rank-deficient source as a positive measure. In
particular no continuation through a competing-hit seam is used.

Use the exact first-physical-defect partition of
Section~\ref{sec:v48-physical-remainder}. If its first type is clearance,
let $j$ be the corresponding flight. Select a responsible nonincident
disk by a fixed ordering of the finitely many candidates, with the
inherited flight and type ordering left unchanged. The responsible
perpendicular foot lies inside the physical segment, by
Lemma~\ref{lem:v49-physical-envelope}. In coordinates based at the
actual terminal collision of this flight, its label is $d\in\mathcal D$.
Set, on the corresponding regular word,
\begin{equation}\label{eq:v61-clearance-map}
 Z_{j,d,R}(x)=Z_d(T_R^{j+1}x),\qquad
 \mathcal J_{j,d,m,R}(x)
   =\det D_{(\alpha,p)}(L_{m,R},Z_{j,d,R})(x).
\end{equation}
The definition of $Z_d$ is \eqref{eq:v49-backward-distance}. The
mark is $j+1$, not $j$, also for the last flight. The first-defect
transition region has
\[
 R\le |Z_{j,d,R}|\le R+s_j,\qquad
 s_j=2\varepsilon q_0^{\min(j,m-1-j)/4},\quad q_0=47/53.
\]
We assume $2\varepsilon<s_0$. The additional candidate selection only
partitions the original positive clearance source; it changes no event.

For a fixed $\eta>0$, split that source according as the absolute
determinant in \eqref{eq:v61-clearance-map}, at its selected witness,
is at least $\eta$ or is less than $\eta$. Denote the resulting exact
roof densities by $b^{\varepsilon,\mathrm{tr},\eta}$ and
$b^{\varepsilon,\mathrm{tan},\eta}$, respectively. Thus
\begin{equation}\label{eq:v61-clearance-positive-split}
 b^{\varepsilon,\mathrm{clr},1}
     =b^{\varepsilon,\mathrm{tr},\eta,1}
                  +b^{\varepsilon,\mathrm{tan},\eta,1},
 \qquad b^{\varepsilon,\mathrm{tr},\eta,1},
        b^{\varepsilon,\mathrm{tan},\eta,1}\ge0.
\end{equation}
The name ``tan'' includes small determinants, not only exact tangencies.

\subsection{Multiplicity of the physical two-coordinate map}
\begin{lemma}[Uniform finite-word two-coordinate multiplicity]
\label{lem:v61-clearance-multiplicity}
There are $C>0$ and $D>1$ such that, for every $m,R,j,d$, the number
of regular preimages of any $(t,z)$ under
$(L_{m,R},Z_{j,d,R})$ at which its determinant is nonzero is at most
$CD^m$, after all physical words and a fixed disjoint angular chart
cover have been included.
\end{lemma}
\begin{proof}
Retain the auxiliary collision graph used in
Lemma~\ref{lem:v44-level-complexity}: contact positions, normals,
unit velocities, flight lengths and the physical first-hit inequalities.
For each fixed center word it has $O(m+1)$ variables and polynomial
sign conditions of fixed degree. The equations $L_{m,R}=t$ and
$Z_{j,d,R}=z$ add only a fixed number of conditions. Indeed $Z_d$
is a scalar product of the image collision position and velocity;
one may use those variables directly, without differentiating a root
or eliminating the intermediate contacts. Rational chart denominators
have a fixed positive sign.

At a regular preimage with nonzero determinant the inverse function
theorem makes that preimage isolated in the fiber. The auxiliary
variables are uniquely determined by the physical initial state and
word, so it gives an isolated component of the same graph fiber.
The coefficient-independent bound \eqref{eq:v44-sign-bound} is
$CD_0^m$ for its number of components. This also bounds the isolated
components when the fiber has other, positive-dimensional components.
There are at most $D_1^m$ physical center words and a fixed number
of charts. Their sum proves the assertion. All constants are uniform
in $R,t,z$ and the chosen mark. Singular or repeated-hit continuations
are not inserted into the graph.
\end{proof}

\begin{theorem}[Height of the transverse physical clearance source]
\label{thm:v61-transverse-clearance}
There are constants $C>0$ and $A>1$ such that, for every $m\ge2$,
$0<2\varepsilon<s_0$, $\eta>0$, and insertions with $|w|\le M$,
\begin{equation}\label{eq:v61-transverse-clearance-height}
 \operatorname*{ess\,sup}_t\sum_{n=1}^{m}\sum_{k\in\Z^2}
 |b^{\varepsilon,\mathrm{tr},\eta,w}_{n,k,m,R}(t)|
                         \le CM A^m\frac{\varepsilon}{\eta}.
\end{equation}
The constants are independent of both thresholds, which may depend
on $m$. Every trajectory, label, first defect and next-collision mark
in this formula is the original one.
\end{theorem}
\begin{proof}
First fix the flight and the candidate disk. On the region where
$|\mathcal J_{j,d,m,R}|\ge\eta$, apply the two-dimensional area
formula to \eqref{eq:v61-clearance-map}. The full collision density
in $(\alpha,p)$ is $1/(4\pi)$; the section restriction multiplies
it by at most $1/c_*$. The density at $(t,z)$ is bounded by
\[
 \frac1{4\pi c_*}\sum_{x:(L_m,Z_{j,d,R})(x)=(t,z),\,
                       |\mathcal J(x)|\ge\eta}|\mathcal J(x)|^{-1}
       \le \frac{CD^m}{c_*\eta}.
\]
One may first exhaust compact regular subsets on which this is an
ordinary finite-to-one change of variables, and then pass by monotone
convergence. Lemma~\ref{lem:v61-clearance-multiplicity} controls the
full multiplicity throughout. The first-defect weight $1-H^{\varepsilon}$, the section event and every exact-label restriction
are all between zero and one, so deleting them here is only a positive
upper comparison on the same physical map.

Integrate $z$ over $[R,R+s_j]\cup[-R-s_j,-R]$, whose total length
is $2s_j$. Summing over candidates costs a fixed factor. Finally
$\sum_{j=0}^{m-1}s_j\le4\varepsilon/(1-q_0^{1/4})$, independently
of $m$. This proves \eqref{eq:v61-transverse-clearance-height} for
the positive source. For fixed $m$, the original events
$\Pi_{n,k,m,R}$ are disjoint. Positive domination proves the sum of
absolute inserted densities in the statement, without multiplying by
the number of exact targets. Fubini and the area formula remove only
Lebesgue-null roof values. No small interval of roofs is discarded.
\end{proof}

\subsection{An exact location of the uncontrolled clearance term}
Use $P=m^2p$, $G=\mathcal L_{m,R}$ and
$\varepsilon(B)=A_0B^{-1/12}$ from
Theorem~\ref{thm:v51-positive-error}. Define the actual finite-count
quantity
\[
 e_{B,m}=\sup_{R,n,k}\|P-G-m^2b^{\varepsilon(B),1}\|_\infty.
\]
For each sufficiently large fixed $B$ its collision limsup is at most
$CB^{-1/192}$. The two preceding theorems and the positive source
identities give the finite inequality
\begin{equation}\label{eq:v61-rank-localized-error}
 \sup_{R,n,k}\|P-G-m^2b^{\varepsilon(B),\mathrm{tan},\eta,1}\|_\infty
 \le e_{B,m}+Cm^2A^m\varepsilon(B)(1+\eta^{-1}).
\end{equation}
Here $A$ is enlarged to work for both geometric bounds. The identity
behind this inequality retains the complete incidence source, transverse
clearance source, and rank-deficient clearance source, with positive
signs. It is not an interchange of the collision and band limits.

In particular a normal-to-normal critical point of a regular word has
$dL_m=0$, and hence $\mathcal J_{j,d,m,R}=0$ for every candidate.
A clearance defect meeting such a point belongs to the retained term,
not to a region on which \eqref{eq:v61-transverse-clearance-height}
can be used. The intrinsic critical jump and all arithmetic phases
remain in the original raw inversion problem.

\paragraph{Relation to the height criterion.}
These results verify two new density-level facts about the Lorentz
source itself: simultaneous inverse-incidence cancellation, and
clearance coarea at a quantified nonzero rank. They are not transfers
from the Markov realization and do not require a monotone scalar
coordinate on every Lorentz word. They use the endpoint action and
coefficient-independent collision-graph complexity already established
for this family. The elementary change-of-variables mechanism is
separate from the dynamical local-limit inputs of \cite{SV,DPZ,DN}.

At a fixed protection width the factor $m^2A^m$ in
\eqref{eq:v61-rank-localized-error} need not be small. Choosing $B$
exponentially in $m$ to offset it would require a reconstruction estimate
not supplied by a fixed-band theorem. Thus neither this inequality nor
the incidence gain alone proves either ordered essential-height limit
in Corollary~\ref{cor:v51-positive-criterion}. A central-scale bound
for the positive exact-label coarea sums, and control of the retained
rank-deficient clearance source, are still required. The unrestricted
pointwise conclusions are not invoked before these steps.
